\section{Introduction}
\label{sec:intro}

Supervised learning needs labeled data, and hand-labeling is the
recurring bottleneck. Programmatic weak supervision attacks the
bottleneck by replacing hand-labels with \emph{labeling functions}
(LFs): small, noisy, user-written programs that vote on an example's
label or abstain \citep{ratner2016data,ratner2017snorkel}. A heuristic
keyword match, a regular expression, a knowledge-base lookup, or the
output of an off-the-shelf classifier each becomes one LF. Many LFs,
individually unreliable, collectively carry signal. The data
programming paradigm \citep{ratner2016data} and the Snorkel system
built on it \citep{ratner2017snorkel} fit a \emph{label model} that
estimates each LF's accuracy and aggregates the LF votes into
probabilistic training labels, which then train a downstream
classifier. The defining promise is gold-free supervision: no
hand-labeled examples are required to fit the label model.

That promise rests on an identifiability claim. The label model is
estimated from LF votes alone, so its parameters (per-LF accuracies,
the class prior) must be \emph{recoverable} from the distribution of
LF votes. When they are not, the fitted label model is one of many
that explain the votes equally well, and the probabilistic training
labels it emits are not trustworthy. The weak-supervision literature
has answered this question in pieces. \citet{dawid1979maximum}
established maximum-likelihood estimation of rater error rates by EM
in the crowdsourcing setting. \citet{ratner2016data} and the triplet
method of \citet{fu2020fast} proved that LF accuracies are
identifiable from agreement moments when the LFs are conditionally
independent given the label. \citet{ratner2019training} extended the
moment approach to structured dependence with a known dependency
graph. Closest to the present setting, \citet{yu2022nplm} let LFs emit
\emph{subsets} of candidate labels, fit a gold-free generative label
model over those subsets, and prove it is identifiable up to label
swapping under mild conditions; our work recovers their candidate-set
view and their label-swap obstruction (our \cref{sec:translation} and
\cref{thm:glass-ceiling}) as the weak-supervision instance of the
reliability-theoretic coarsening framework, and adds the
dependent-case gold-set sample-complexity analysis they do not treat.
What this body of work has not assembled is a single inferential frame
that classifies \emph{which} LF ensembles are identifiable, names the
auxiliary information that rescues the rest, and prices it.

\subsection{The bridge to masked-cause inference}

We observe that label-model estimation is an instance of the
masked-cause series-system identifiability problem from reliability
statistics \citep{towell2026masked}. In that framework, $m$
components in series produce a system that fails when any one fails;
the failed component is masked, but a \emph{candidate set} containing
it is observed. Identifiability of component-level parameters hinges
on coarsening-at-random conditions, the C1--C2--C3 conditions of
\citet{heitjan1991ignorability} and \citet{gill1997coarsening}, and a
\emph{singleton} candidate set (an observation where the cause is
known exactly) restores identifiability when broader coarsening is
confounded.

The translation to weak supervision is direct. The true label $Y$ is
the latent cause. The set of labels consistent with the LF votes for
an example is the candidate set. An LF that abstains contributes
nothing and leaves the candidate set full (a non-informative
coarsening); an LF that votes confidently narrows the candidate set;
a gold-labeled (hand-annotated) example pins the candidate set to a
singleton. The C1--C2--C3 conditions become a classification of LF
ensembles, and the masked-cause identifiability theorems give the
precise conditions under which a gold-free pipeline is identifiable.
The present paper is one application of the masked-cause framework
of \citet{towell2026masked}.\footnote{Sibling applications in the
framework series address scRNA-seq zero inflation, spatial
transcriptomics, differential privacy, and EHR phenotyping
\citep{towell2026scrnacoarsening,towell2026spatialcoarsening,
towell2026dpcoarsening,towell2026phenotypecoarsening}.}

\subsection{Contributions}

We contribute:

\begin{itemize}
\item \textbf{A bridge from programmatic weak supervision to
  masked-cause inference} (\cref{sec:translation}), expressing the
  label model of data programming and Snorkel as a coarsening problem
  and casting the C1--C2--C3 conditions as a classification of LF
  ensembles.

\item \textbf{A glass-ceiling theorem} (T1, \cref{sec:identifiability}):
  without gold labels and without a structural assumption on LF
  dependence, the LF accuracies and class prior are non-identifiable.
  We exhibit the obstruction explicitly as an accuracy-complement
  symmetry, broken only by an assumption such as ``LFs are better than
  random'' or conditional independence.

\item \textbf{An identifiability theorem under conditional
  independence} (T2, \cref{sec:identifiability}): with conditionally
  independent LFs and at least three of non-degenerate accuracy, the
  LF accuracies and class prior are identifiable from the second- and
  third-order agreement moments. This re-derives the
  Dawid--Skene / data-programming / triplet-method result as a
  masked-cause theorem.

\item \textbf{An agreement-consistency theorem} (T3,
  \cref{sec:identifiability}): the fitted label model reproduces the
  empirical pairwise LF agreement rates exactly at an interior MLE.
  This is the label-model analogue of the cell-total consistency
  result of \citet{towell2026scrnacoarsening}.

\item \textbf{A gold-set sample-complexity theorem} (T4,
  \cref{sec:methodology}): when LFs violate conditional independence,
  with the dependence inducing a parameter-degeneracy dimension $r$ in
  the agreement structure, the number of gold-labeled examples needed
  to restore the label model to total (L2) accuracy is
  $\Theta(r / \mathrm{gap}^2)$ in that degeneracy dimension and the
  accuracy margin, with matching upper and minimax lower bounds; the
  rate is loss-dependent, falling to $\Theta(\log r / \mathrm{gap}^2)$
  under a per-direction criterion. Singletons restore the rank that LF
  dependence destroys.

\item \textbf{Simulation validation} (\cref{sec:validation}). A
  base-R study confirms the glass-ceiling symmetry, gold-free recovery
  under conditional independence, agreement consistency to numerical
  precision, and the gold-set scaling under injected LF dependence.
\end{itemize}

The paper's message: gold-free weak supervision is not magic, it is a
coarsening problem with a precise identifiability boundary. The
C1--C2--C3 conditions tell a practitioner whether their LF ensemble
is on the identifiable side of that boundary, and the
sample-complexity theorem tells them how much gold data buys back
identifiability when it is not.
